%\input ../util/bookmacros
%\input ../util/docparams

\eject
\section{References}

The works cited in this book, together with the closest kindred spirits of
its educational point of view.

\def\refentry#1{\smallskip{\hangindent=2em\hangafter=1\noindent #1\par}}

\medskip
\refentry{Ayres, F., and E. Mendelson, {\it Schaum's Outline of
Calculus\/}, 6th ed., McGraw-Hill, 2013. (The drill companion recommended
in the preface and in the appendix ``How to Study from This Book.'')}

\refentry{Bjork, E. L., and R. A. Bjork, ``Making things hard on yourself,
but in a good way: Creating desirable difficulties to enhance learning,''
in M. A. Gernsbacher et al.\ (eds.), {\it Psychology and the Real World:
Essays Illustrating Fundamental Contributions to Society\/}, Worth
Publishers, 2011, 56--64. (The ``spacing'' and ``desirable difficulties''
findings behind the study schedule of the final appendix.)}

\refentry{Bressoud, D., {\it Calculus Reordered: A History of the Big
Ideas\/}, Princeton University Press, 2019. (A kindred argument that the
standard order of the calculus curriculum is a historical accident, not a
learning path.)}

\refentry{Cuoco, A., E.~P. Goldenberg, and J. Mark, ``Habits of Mind: An
Organizing Principle for Mathematics Curricula,'' {\it Journal of
Mathematical Behavior\/} 15 (1996), 375--402. (The thesis that recurring
ways of thinking, not topics, should organize a curriculum -- the nearest
published relative of this book's four strategic moves.)}

\refentry{Ebbinghaus, H., {\it \"Uber das Ged\"achtnis\/}, 1885; in
English: {\it Memory: A Contribution to Experimental Psychology\/},
Teachers College, Columbia University, 1913. (The original measurements
of forgetting.)}

\refentry{Freudenthal, H., {\it Revisiting Mathematics Education: China
Lectures\/}, Kluwer, 1991. (``Guided reinvention'': the student should
experience mathematics as knowledge they could have made their own -- this
book's ``player'' principle, stated a generation earlier.)}

\refentry{Harel, G., ``Intellectual Need,'' in K.~R. Leatham (ed.), {\it
Vital Directions for Mathematics Education Research\/}, Springer, 2013,
119--151. (The research formulation of ``a tool should answer a question
the student already has'' -- this book's Problem Driven Development.)}

\refentry{P\'olya, G., {\it How to Solve It\/}, Princeton University
Press, 1945. (The tradition of teachable strategic moves to which this
book's four moves belong.)}

\refentry{Roediger, H.~L., and J.~D. Karpicke, ``Test-Enhanced Learning:
Taking Memory Tests Improves Long-Term Retention,'' {\it Psychological
Science\/} 17 (2006), 249--255. (Retrieval beats rereading -- the first
finding of the final appendix.)}

\refentry{Skemp, R.~R., ``Relational Understanding and Instrumental
Understanding,'' {\it Mathematics Teaching\/} 77 (1976), 20--26. (The
classic statement of the difference between understanding a subject and
executing its procedures.)}

\medskip
